\cotefile{93}
\scanpage{32}
% archive sheet number: 31
\mspage{6}

Let $y$ be in $Y$. Then $\theta(y)$ is in the norm closed convex hull of $\{t(y):t\in T\}$ by [3, Th. 13, p. 422]. Therefore, given any $\varepsilon>0$ there exist $t_1,\ldots,t_n$ in $T$ and $\alpha_i\geqq 0$ $(i=1,\ldots,n)$ such that $\sum\alpha_i=1$ and
\[
\lVert\theta(y)=\sum\alpha_i t_i(y)\rVert<\varepsilon/(M+1),
\]
where $M=\sup\{\lVert t\rVert:t\in T\}$ (which is finite by the uniform boundedness principle.) Then if $t\geqq t_1\unsure{\mathbin{\mathrm{o}}|\circ}\ldots\mathbin{\mathrm{o}}t_n$,
\[
\begin{aligned}
\lVert\theta(y)-t(y)\rVert&=\lVert t[\theta(y)-\sum\alpha_i t_i(y)]\rVert\\
&\leqq M\varepsilon/(M+1)<\varepsilon,
\end{aligned}
\]
so $\displaystyle\lim_{t\in T}\lVert\theta(y)-t(y)\rVert=0$. \hfill Q.E.D.

In order to apply the theorem to the case of $T=\{t_\pi:\pi\in\unsure{\mathcal{P}|P}\}$ acting on $\mathrm{sba}(\unsure{\mathcal{A}|A},X)$, we must show that for every $\mu$ in $\mathrm{sba}(\mathcal{A},X)$ the set
\[
\{t_\pi\mu:\pi\in\mathcal{P}\}
\]
is conditionally weakly compact in $\mathrm{sba}(\mathcal{A},X)$. This follows immediately from Proposition 1 below. The theorem then implies that $\mathrm{sba}(\mathcal{A},X)$ can be decomposed into the direct sum of the countably additive members of $\mathrm{sba}(\mathcal{A},X)$ and the purely finitely additive members according to the following definition.

\scanpage{33}
% archive sheet number: 32
% slanted
\margin{hyp. de \unsure{l'induction|induction}}
$3^\circ)$ L'hom can
\[
\Omega_X|Y\longrightarrow\underline{\mathrm{Hom}}(\unsure{P|\rho}_{Y/X},\Omega_Y)
\]
est injectif [une orientation de $X$ en $y\in Y$ est connue quand on sait comment \inserted{au vois. de $y$} une rive \inserted{de $Y$\deleted{\unsure{ en $y$}}} détermine une orientation de $Y$]

$4^\circ)$ Pour tte $\omega\in\Gamma(X)$, l'hom correspondant
\[
P_{Y/X}\longrightarrow\Omega_Y
\]
% slanted
\margin{hyp. de non trivialité}
est injectif [deux rives différentes définissent des orientations différentes de $Y$]

$5^\circ)$ Pour tout $X\in\mathrm{Ob}\,\unsure{V|\mathcal{V}}$, $\Omega(X)\to X$ est couvrant (il existe localement des orientations \unsure{$g$|$?$} de $X$; \inserted{ou encore : $\forall n\geq 1$, $\Omega(\mathbb{R}^n)$ est non vide})

% slanted
\margin{\deleted{\unsure{hyp.}}principe de \deleted{\unsure{pseudo}}l'identité}
$6^\circ)$ Pour tt $X\in\mathrm{Ob}\,V$, $\Omega(X)$ interprété comme espace étalé sur $X$ est séparé sur $X$, ou encore : une orientation \unsure{de l'un|de l'espace|d'une} $X\in\mathrm{Ob}\,V$ est connue en un pt quand on la connait en des pts arbitrairement voisins. (Il suffit de l'exiger pour $\mathbb{R}^n$, $n\geq 1$, en fait, il sera vrai à posteriori qu'une orientation de $\mathbb{R}^n$, \inserted{ou de tt $X$ connexe,} est connue quand on la connaît en un seul point\ldots)

\scanpage{34}
% archive sheet number: 33
\mspage{5}

(iii) \uline{The} \uline{space} $Y$ \uline{can} \uline{be} \uline{written} \uline{as} \uline{a} \uline{direct} \uline{sum}
\[
Y=Y_1\oplus Y_0,
\]
\uline{where}
\[
\begin{aligned}
Y_1&=\{y\in Y:y\text{ is a }T\text{-fixed point}\}\\
&=\theta(Y)
\end{aligned}
\]
\uline{and}
\[
\begin{aligned}
Y_0&=\{y\in Y:0\in\unsure{\mathcal{O}|O}(y)\}\\
&=\text{kernel of }\theta.
\end{aligned}
\]

PROOF. The assumption that each $\mathcal{O}(y)$ is weakly compact insures that the weak operator closure $S$ of $T$ in $B(Y)$ is compact in the weak operator topology. Moreover, $\mathcal{O}(y)=\{s(y):s\in S\}$ for every $y$ in $Y$. By the lemma, $T$ converges in the weak operator topology to a zero $\theta$ for $S$. Since $t\unsure{\mathbin{\mathrm{o}}|\circ}\theta=\theta$ for all $t$ in $T$, $\deleted{\unsure{\mathcal{O}|O}}\handwritten{\theta}(y)$ is $T$-fixed for any $y$. If $z$ is any $T$-fixed point in $\deleted{\theta}\handwritten{\unsure{\mathcal{O}|O}}(y)$, then $z=s(y)$ for some $s$ in $S$ and
\[
z=\theta(z)=\theta(s(y))=\theta(y).
\]
\unsure{;|,}

Hence (i) and the first part of statement (ii) hold; (iii) follows trivially. It remains only to show that $T$ converges to $\theta$ in the strong operator topology.

\scanpage{36}
% archive sheet number: 35
\[
g=\exp f\qquad f'=\varphi=\psi^{1/2}
\]
\[
g'=(\exp f)f'\ \deleted{\unsure{=(\exp f)\varphi}}
\]
\[
g''=(\exp f)f'^{\,2}+(\exp f)f''=(\exp f)(f'^{\,2}+f'')
\]
\[
\begin{aligned}
g'''&=(\exp f)f'(f'^{\,2}+f'')+(\exp f)(2f'f''+f''')\\
&=(\exp f)(f'^{\,3}+3f'f''+f''')
\end{aligned}
\]
\[
\frac{g'''}{g'}=f'^{\,2}+3f''+\frac{f'''}{f'}
\]
\[
\frac{g''}{g'}=f'+\frac{f''}{f'}\qquad
\left(\frac{g''}{g'}\right)^2=f'^{\,2}+\left(\frac{f''}{f'}\right)^2+2f''
\]
\[
\begin{aligned}
-2\frac{g'''}{g'}+3\left(\frac{\unsure{f|g}''}{g'}\right)^2
&=\underline{-2f'^{\,2}}-\deleted{6f''}-2\frac{f'''}{f'}\\
&\quad+\underline{3f'^{\,2}}+3\left(\frac{f''}{f'}\right)^2+\deleted{6f''}\\
&=-\deleted{\unsure{3}}2\frac{f'''}{f'}+3\left(\frac{f''}{f'}\right)^2+\deleted{\unsure{3}}f'^{\,2}
\end{aligned}
\]
\[
G(t,\exp f)=G(t,f)+3(df)^2.
\]
\[
\begin{aligned}
&\ \ \Vert\\[-2pt]
-C_t+C_{\exp f}&=-C_t+C_f+3(df)^2
\end{aligned}
\]
\[
\boxed{C_{\exp f}=C_f+\deleted{\unsure{3}}df^2}
\]
\[
C_{\exp f}=C_f+\deleted{\unsure{3}}df^2=C_f+\omega^{\otimes 2}
\]
\[
\begin{aligned}
\deleted{\unsure{g}}\varphi&=\exp \unsure{t|b}\\
\varphi'&=\exp t\\
\varphi''&=\exp t\\
\varphi'''&=\exp t
\end{aligned}
\qquad
\left(-2\frac{\varphi'''}{\varphi'}+3\left(\frac{\varphi''}{\varphi'}\right)^2\right)dt^2=dt^2
\]
\[
\unsure{C|c}_{\!f}\quad \unsure{+\infty|+\omega}
\]
% Opposite orientation, beginning at the lower edge of the scan.
\[
\deleted{\unsure{G|\ell}}G(z,f)\deleted{\unsure{df}}=
\left(-\unsure{2}\frac{f'''}{f'}+3\left(\frac{f''}{f'}\right)^2\right)dz^2
\]
\margin{\unsure{invariance|invariante} \unsure{par|pour} $\pi$ \deleted{\unsure{paraboliques|projectifs}} ``\unsure{paraboliques|projectives|hyperboliques}''}
\[
\frac{d\deleted{\unsure{z}}f}{d\deleted{\unsure{f}}z}\quad \text{forme inv. par }\pi
\]
\[
G(f,z)\deleted{df^2}=0
\]
\[
\frac{d^3z}{df^3}\ \unsure{\sim|=}\ 0
\]
\[
G(f,z)=0
\]
\[
\begin{aligned}
G(z,f)&=\\
\unsure{\gamma|r}\mathbin{.}G(z,f)&\\
\Vert&\\
G(\gamma\mathbin{.}z,\gamma\mathbin{.}f)&\\
\Vert&\\
G(\gamma z,f)&\\
\Vert&\\
G(z,f)&
\end{aligned}
\]

\scanpage{37}
% archive sheet number: 36
$\overline{\underline{\mathrm{V}}}$ Programme de \unsure{la|le} Théorie

(Cristaux)

\scanpage{38}
% archive sheet number: 37
% slanted
\margin{\handwritten{\unsure{Nouveau|Nouveaux} Chapitre}}
\handwritten{1}) Programme de la théorie de Dieudonné sur une base \inserted{$S$} où $p$ est localement nilpotent.

a) \uline{Foncteur de Dieudonné}. On \deleted{va} définiraun foncteur
\margin{\handwritten{$\mathbb{D}^*$}}
\[
\underline{D}^{*}:\quad \mathrm{BT}(S)^o\handwritten{\longrightarrow}\mathrm{Crisloclib}(S)\quad,
\]
compatible aux images inverses (donc foncteur cartésien sur la catégorie fibrée des groupes de BT sur des bases où $p$ est loc\deleted{\unsure{,|;}}\handwritten{\unsure{.|;}} nilp.). Si $S_{\unsure{o|0}}=\mathrm{Var}(p)$ on a une équivalence de catégories (cf.\qquad plus haut)
\[
\mathrm{Crisloclib}(S)\handwritten{\approxeq}\mathrm{Crisloclib}(S_o)
\]
de sorte que la donnée de $\underline{D}^{*}$ équivaut à la donnée de
\[
\underline{D}^{*}:\quad\mathrm{BT}(S_o)^o\handwritten{\longrightarrow}\mathrm{Crisloclib}(S)_o\quad,
\]
i.e. on est ramené \deleted{à la car. $p$} au cas $S$ de car. $p$. \deleted{Pour connaitre, da}

Dans le cas général, $\underline{D}^{*}(G)$ ne dépend que de \deleted{$\unsure{\underline{D}^*|D^*}$} $G_o=G\mathbin{\mathrm{x}}_S S_o$.

\deleted{Utilis} Plaçons nous sur $\deleted{\unsure{X}}S_o$ de car. $p$. Utilisant $\underline{f}_G$ et $\underline{v}_G$ et la compatibilité de $\underline{D}^{*}$ aux images inverses, on voit que \deleted{le foncteur} $\underline{D}^{*}(G_o)$ \deleted{peut}\inserted{$=M$} est naturellement muni de morphismes $F_M,V_M$ \deleted{satisfaisant}
\diagram{\begin{tikzcd}[ampersand replacement=\&]
M^{(p)}=\underline{f}_S^{*}(M) \arrow[r,shift left,"\handwritten{F_M}"] \& M \arrow[l,shift left,"\handwritten{V_M}"]
\end{tikzcd}}
satisfaisant
\[
V_M F_M=p\operatorname{id}_M,\quad F_M V_M=p\operatorname{id}_{\deleted{\unsure{m}}M^{(p)}}\quad.
\]
\deleted{\unsure{On}}On appelle un tel triple $(M,F_M,V_M)$ un (F-V-\uline{cristal}\inserted{\handwritten{ ou \uline{cristal de Dieudonné})}}, et on a donc un foncteur
\[
\underline{D}^{*}:\quad\mathrm{BT}(S_o)\handwritten{\longrightarrow}\text{F-V-Cris}(S_o)\quad .
\]
C'est ce foncteur qui mérite le nom de \uline{foncteur de Dieudonné} . On \deleted{\unsure{pura|poura} prouver que} voit que sa formation commute à tout changement de\handwritten{\inserted{ }}base. De plus, on prouve que pour $S_o$ spectre d'un corps parfait, il coïncide \deleted{avec} à isomorphisme canonique près avec le foncteur de Dieudonné habituel, qui est alors une équivalence de catégories, de sorte que dans ce cas la

\scanpage{39}
% archive sheet number: 38
connaissance de $\underline{D}^{*}(G_{\unsure{o|0}})$ permet de réconstituer $G_o$. Ainsi, \deleted{\unsure{a}}dans le cas général, la connaissance de $\underline{D}^{*}(G_o)=M$ comme F-V-cristal sur $S_o$ permet de retrouver les fibres de $G_o$ en des corps parfaits au dessus de $S_o$ ; cela justifie donc le sentiment que pour l'essentiel, on a saisi avec le cristal $M=\underline{D}^{*}(G_o)$ la \deleted{\unsure{f}}famille des $G_{o\bar{s}}$. Une question importante \inserted{\handwritten{en ce sens}} que je n'ai pas résolue est :

\uline{Problème} 1 : Le foncteur de Dieudonné sur une base $S_o$ de car. $p$ est-il pleinement fidèle ? \deleted{Quel est}\deleted{\inserted{\unsure{le|la}}} \deleted{\unsure{s}}Son image essentielle \deleted{\unsure{e|S}}est-elle formée des F-V-cristaux ``admissibles'' définis plus bas ?

Pour résoudre cette question, il me semble qu'il faudrait arriver à formuler une théorie de Dieudonné s\deleted{\unsure{r|u}}ur \deleted{des bases de car $p$} $S_o$ pour des schémas en \inserted{$p$-}groupes finis localement libres sur $S_o$ quelconques, ou du moins pour ceux plats sur un $\handwritten{\Lambda}_n$, en trouvant une équivalence entre \deleted{cette} la catégorie de\deleted{s} ces groupes sur $S_o$, et une catégorie de F-V-cristaux sur $S_o$ \uline{muni de structures supplémentaires}\inserted{\handwritten{$^{(1)}$}}\margin{\handwritten{sans doute inutile si $S_o$ est parfait}} que je n'ai pas réussi à dégager encore, et qui très probablement impliqueront la donnée d'une \uline{filtration} sur un objet \inserted{convenable} d'une catégorie dérivée \ldots

\scanpage{40}
% archive sheet number: 39
\uline{Problè\deleted{\unsure{p|m}}me} 2. Développer une telle théorie de Dieudonné pour des $p$-Groupes finis localement libres sur une base $S_{\unsure{o|0}}$ de car.$p$ \ $0$. Si $S_o$ est parfait établir une \inserted{anti-}équivalence entre la catégorie des \deleted{\unsure{G|g}}groupes finis localement libres sur $S_o$ qui sont \uline{plats sur} $\handwritten{\Lambda}_n=\underline{Z}/p_{\deleted{n}}^n\underline{Z}$ (cf.\qquad), et la catégorie des $\underline{W}_n(\underline{O}_S)$-Modules localement libres $M$, muni de $F_M$ et $V_M$ \deleted{comme d'habi} satisfaisant les conditions habituelles.

Une solution affirmative à ce dernier problème devrait évidemment impliquer une solution affirmative au problème 1 dans le cas où $S_o$ est un schéma \uline{parfait}.

\scanpage{41}
% archive sheet number: 40
b) \uline{Filtration}s\deleted{\unsure{ass}} \uline{associées aux cristaux de Dieudonné}. Revenons au cas d'un $S$ général (sur lequel $p$ soit loc. nilp.). On défini\deleted{\unsure{t}}ra alors, pour tout $G\handwritten{\in}\mathrm{Ob}\,\mathrm{BT}(S)$, une filtration \deleted{en deux crans} d\deleted{e}\unsure{u|e} Mo\deleted{\unsure{d|D}}\unsure{D|d}ule localement libre $\underline{D}^{\handwritten{*}}(G)_S$ sur $S$ par un sous-Module localement facteur direct (donc localement libre de t.f.) $\inserted{\handwritten{\underline{\mathrm{Fil}}^1=}}\underline{\handwritten{\omega}}_G$, donnant lieu à une suite exacte
\[
0\handwritten{\longrightarrow}\underline{\handwritten{\omega}}_G\handwritten{\longrightarrow}\underline{D}^{\handwritten{*}}(G)\handwritten{\longrightarrow}\underline{t}_{G^{\handwritten{*}}}\handwritten{\longrightarrow}0\quad,
\tag*{\handwritten{($\unsure{*|\delta}$)}}
\]
où en fait $\underline{\handwritten{\omega}}_G$ est le module des différentielles invariantes sur le groupe formel $\bar{G}\handwritten{\subset}G$ associé à $G$ (qui sera construit plus bas), et $\underline{t}_G=\underline{\handwritten{\omega}}_{G^{\handwritten{*}}}^{\handwritten{\vee}}$ est l'Algèbre de Lie associée au gro\deleted{\unsure{u|v}}upe de BT dual $G^{\handwritten{*}}$ de $G$ (définition à reporter aux généralités sur les goupes de BT). Cette \deleted{filtr} suite exacte est fonctorielle \inserted{en $G$} \deleted{et} compatible aux images inverses par un $S'\handwritten{\longrightarrow}S$.

\margin{\handwritten{(*) Sans doute inutile si $S_{\unsure{o|0}}$ \unsure{parfait|parfaite}.}}

\scanpage{42}
% archive sheet number: 41
Appelons F-V-\uline{cri\deleted{\unsure{t}}s\deleted{\unsure{s}}tal} \uline{filtré} \inserted{\uline{sur} $S$} (par abus de langage) un F-V-cristal en Mod. loc. lib. sur $S$ muni d\deleted{'}e \deleted{une filtr.} la donnée d'une filtration de $M_S$ par une sous-module localement facteur direct. On trouve donc en fait un foncteur
\[
\mathrm{BT}(S)\handwritten{\longrightarrow}\text{F-V-cris fil }(S)\quad ,
\]
compatible aux images inverses.

\inserted{\uline{Remarque}.} La filtration $(*)$ doit être considérée comme \deleted{\unsure{e}}étant l'é\deleted{\unsure{q}}quivalent d'une filtration ``de Hodge'' sur une cohomo\deleted{\unsure{o}}logie de De Rham \inserted{relative} en dimension 1. \deleted{Pour expliciter} \deleted{\unsure{c}}Ce point\deleted{, partons par exemple d'un schéma abélien $A$ sur $S$, et prenons $G=A(\infty)=\lim A(n)$. On peut alors} \inserted{sera explicité plus bas.}

Soit maintenant $S'$ un épaississement à p.div. de $S$. On se propose de trouver tous les prolongements possibles (à isom. près) d\deleted{\unsure{a}}u groupe de BT $G$ donné sur $S$ en un groupe de BT $G'$ sur $S'$. On \deleted{trouve} prouvera que si à tout tel $G'$, on associe la filtration correspondante de $D(G\handwritten{'})_{S'}=\underline{D}^{\handwritten{*}}(G)_{S'}=\underline{D}^{\handwritten{*}}(G_{\unsure{o|0}})_{S'}$ (cf.\qquad pour le fait que $S'$ peut être envisagé comme un objet du site cristallin absolu de $S_o$), qui prolonge la filtration qu'on a sur $\underline{D}^{\handwritten{*}}(G)_S=\underline{D}^{\handwritten{*}}(G_o)_S=\underline{D}^{\handwritten{*}}(G')\handwritten{|}S$, on trouve une bijection entre l'ensemble des classes d'isomorphie cherchés et l'ensemble des filtrations qui prolongent. De façon plus précise, on a :\deleted{une équivalence de catégories}

\uline{Théorème de déformation pour les groupes de} BT : Soient $S$ un schéma où $p$ est loc. nilp. , $S'$ un épaississement à p.div. de $S$, Considérons le foncteur canonique
\[
\begin{aligned}
\mathrm{BT}(S')\handwritten{\longrightarrow}\quad&
\text{catégorie des \deleted{triples}\inserted{paires} d'un groupe de BT $G$ sur $S$, et}\\
&\text{d'un\deleted{e} sous-module localement facteur direct de}\\
&\underline{D}^{\handwritten{*}}(G)_{S'}\quad\text{qui prolonge le sous-module }\mathrm{Fil}^{\handwritten{1}}\underline{D}^{\handwritten{*}}(G)_S=\underline{\handwritten{\omega}}_G\\
&\text{de }\underline{D}^{\handwritten{*}}(G)_S\ .
\end{aligned}
\]
Ce foncteur est une \uline{équivalence de catégories}. \inserted{\handwritten{\uline{si} $S'$ \uline{est à p. div. nil} ou si on se borne à des $G$, \unsure{$G'$} qui sont infinitésimaux, ou à \unsure{duaux} infinitésimaux.}}

\inserted{\handwritten{Remarques}} Si l\deleted{\unsure{a}}e \deleted{questio} problème 1 avait une réponse affirmative, on en conclurait une description de la catégorie des groupes de BT sur $S$ en termes de cristaux de Dieudonné filtrés sur $S$, en appliquant le théorème de

\scanpage{43}
% archive sheet number: 42
déformation au\deleted{x} \deleted{\unsure{r}}cas du couple $(S_{\unsure{o|0}},S)$, \deleted{en} notant que l'idéal $p\underline{O}_S$ est bien un idéal à puissances divisées. On notera d'ailleurs que le théorème de déformations dans le cas \deleted{envisagé} général $(S,S')$ e\deleted{\unsure{n}}nvisagé dans son énoncé \deleted{\unsure{e}}\unsure{est|ést} en fait équivalent au cas particulier d'un couple $(S_o,S)$, comme il est formel v\deleted{\unsure{l}}ia $\mathrm{Crisloclib}(S)\handwritten{\xrightarrow{\approx}}\mathrm{Crisloclib}(S_o)$.

\inserted{\handwritten{\uline{Problème} 3\quad Trouver variante du th. de déformation pr les $p$-Groupes}}
% slanted
\margin{\handwritten{finis loc. libres (j'ai un candidat).}}

c) F-V-\uline{cristaux admissibles en car}. $p$. Revenons au cas d'un schéma de car. $p$ . Le cristal de Dieudonné $D^{\handwritten{*}}(G_o)$ étant muni d'une filtration canonique (de $D^{\handwritten{*}}(G_o)_{S_o}$), on pourrait craindre que le problème 1 est mal posé, et qu'il faudrait envisager le foncteur comme étant à valeurs dans des cristaux de Dieudonné filtrés. En fait, dans le cas envisagé ici, il se trouve que la filtration est uniquement déterminée en \deleted{y}termes de la structure de cristal de Dieudonné, comme on va l'expliquer maintenant.

Considérons le morphisme de topos annelés (qui aura été défini déjà précédemment, dans le par. d'exemples de cristaux, pour les relations entre cristaux et vecteurs de Witt)
\diagram{\begin{tikzcd}[ampersand replacement=\&]
\phi:\quad S_{o\,\mathrm{c}\deleted{\unsure{r}}\mathrm{ris}/\underline{F}_p}
 \arrow[r] \& S_o \arrow[l,bend right=38,"\handwritten{\varepsilon}"']\quad.
\end{tikzcd}}
Soit $M=D^{\handwritten{*}}(G_o)$, et soit $M_o$ sa restriction au site cristallin rel. à $\underline{F}_p$. Il se trouve que l\deleted{a\unsure{s|{}} suite} d'après $FV=p$, $VF=p$, les composé\deleted{\unsure{z}}s dans la suite
\diagram{\begin{tikzcd}[ampersand replacement=\&]
\deleted{\unsure{O|M}}M_o^{(p)}\arrow[r,"\handwritten{F_{M_o}}"] \&
M_o\arrow[r,"\handwritten{V_{M_o}}"] \&
M_o^{(p)}\arrow[r,"\handwritten{F_{M_o}}"] \& M_o
\end{tikzcd}}
sont nuls; en fait, on prouvera mieux, à savoir que \inserted{\handwritten{pour tout cristal de Dieud. $M$,}} la suite précédente est exacte (terme à terme sur tout épaississement à p.div. de car.$p$ de $S_o$) et on obtient ainsi des sous-cristaux en modules loc. libres
\[
\operatorname{Ker} V_{M_o}=\operatorname{Im} F_{M_o}\handwritten{\subset}M_o,\quad
\operatorname{Ker} F_{M_o}=\operatorname{Im} V_{M_o}\handwritten{\subset}M_o^{(p)}\quad .
\]
Ceci dit, \deleted{\unsure{le|de} sous-Module $\mathrm{Fil}^1$} notons que $M_{S_o}=\handwritten{\varepsilon^*}(M_o)$, donc
\[
\phi^{\handwritten{*}}(M_{S_o})\handwritten{\simeq} M_o^{(p)}
\]
donc on a
\[
\phi^{\handwritten{*}}(\deleted{\unsure{E|F}}\mathrm{Fil}^1(M_{S_o}))\handwritten{\subset} M_o^{(p)}\quad .
\]
Ceci posé, on montre qu'on a

\scanpage{44}
% archive sheet number: 43
% The circled text at the head is called in after « Fil¹ est unique et ».
\[
\phi^{\handwritten{*}}(\mathrm{Fil}^1(M_{S_{\unsure{o|0}}}))=
\operatorname{Ker} F_{M_o}\ (=\operatorname{Im}V_{M_o})\handwritten{\subset}M_o^{(p)}\quad .
\]
D'a\deleted{\unsure{n}}utre part, on montre que $\phi$ a des propriétés de fidélité telles que la relation précédente détermine sans ambiguité le sous-Module localement facteur direct $\mathrm{Fil}^1$ de $M_{S_o}$.

On dira qu'un \inserted{F-V-}c\deleted{\unsure{ri}}ristal en modules loc. lib. $M_o$ sur le site cristallin de $S_o$ au des\deleted{\unsure{s}}sus de $\underline{F}_p$ est \uline{admissible} s'il existe un sous-module l\deleted{\unsure{a}}oc. \deleted{facteu} facteur direct $\mathrm{Fil}^1$ de $M_{oS}$ tel que l'on ait la relation précédente\unsure{;|,} \deleted{On} \inserted{$\mathrm{Fil}^1$ est unique et \deleted{déterm}\unsure{$\circ$|o}\inserted{détermine une filtration qu'on appellera la \uline{filt}. \uline{canonique} . \handwritten{On}}} dit de même qu'un \deleted{\unsure{Fil}}cristal de Dieudonné \inserted{$M$} sur $S_o$ est admissible si \deleted{\unsure{l}}sa restriction $M_o$ \deleted{sur} au site cristallin rel. à $\underline{F}_p$ est admissible. On voit donc que si $M$ est le cristal de Dieudonné associé à un groupe de BT sur $S_o$, $M$ est admissible, et la filtration de $M_{S_o}$ envisagée dans b) est la filtra- canonique définie précédemment.

% slanted
\margin{\handwritten{d) Passage à la limite : \unsure{prendre|pour des|formule} anneau $A$ \unsure{séparé} et complet pour top. \uline{\unsure{$p$-adique}}}}
% slanted
\margin{\handwritten{e) Théorie \unsure{:|à} isogénie près \uline{et} th. de Tate}}
% slanted
\margin{\handwritten{f) Dualité}}

\handwritten{2}) \uline{Théorie de déformations pour les schémas abéliens}.

\handwritten{a)} Dans cette théorie, il n'y a plus de nombre premier privilégié, par contre il faut travailler avec le site cristallin \uline{nilpotent}. On trouve un foncteur \inserted{contravariant} de la catégorie des schémas abéliens sur le schéma $S$ vers les \deleted{sites} cristaux sur le site cristallin nilpotent de $S$
\[
\text{Schémab}(S)\handwritten{\longrightarrow}\mathrm{Crisloclib}_{\mathrm{nilp}}(S)
\]
qui peut être défini par exemple comme
\[
\underline{\handwritten{\mathbb{D}}}^{*}:\ A\handwritten{\longmapsto}R^1 f_{A\mathrm{cris}*}(\deleted{A}\handwritten{O_{A_{\mathrm{cris}}}})
\tag*{\handwritten{(*)}}
\]
où on considère
\[
f_{\mathrm{cris}}:\quad A_{\mathrm{crisnilp}}\handwritten{\longrightarrow}S_{\mathrm{crisnilp}}
\]
induit par $S$. Par \deleted{\unsure{l}}une va\deleted{\unsure{i}}riante du théorème de changement de base, on voit que
\[
\mathbb{D}^*(A)_S=\deleted{\unsure{R^1 f_{A\mathrm{DR}*} A}}\underline{H}_{\mathrm{DR}}^1(A/S)\overset{\mathrm{dfn}}{=}R^1 f_{A*}(\handwritten{\Omega}_{A/S}^{\bullet})\quad .
\]
Pour calculer $\mathbb{D}^{\handwritten{*}}(A)_{S'}$, $S'$ épaississement à p.div. loc. nilpotentes de $S$, on aura \deleted{\unsure{p}}pour tout prolongement \deleted{canonique} de $A$ en un schéma abélien $A'$ sur $S'$

\scanpage{45}
% archive sheet number: 44
\[
\mathbb{D}^{\handwritten{*}}(A)_{S'}=\deleted{R^1 f}\underline{H}_{\mathrm{DR}}^1(A'/S')\quad\deleted{\unsure{,|;}}.
\]
\deleted{\unsure{o|6}}On notera que d'après un théorème connu (cf. li\deleted{\unsure{b}}vre de Mumford Geometric Invariant theory) il existe toujours un $A'$ pour $S$ affine; le fait que \deleted{\unsure{il|y}}\inserted{$H^1(A'/S')$} ne dépende, à isomorphisme près, pas du prolongement choisi de $A$, peut être considéré ici comme conséquence du tapis de la cohomologie cristalline.

En fait, utilisant la filtration de \deleted{De Rham}\inserted{Hodge} sur la cohomologie de De Rham, on trouve une filtration sur $\mathbb{D}^{\handwritten{*}}(A)_S$, savoir
\[
0\handwritten{\longrightarrow}R^{\unsure{o|0}} f_{\handwritten{*}}(\underline{\handwritten{\Omega}}^1_{A/S})\handwritten{\longrightarrow}\underline{H}_{\mathrm{DR}}^1(A/S)\handwritten{\longrightarrow}R^1 f_*(\underline{O}_S)\handwritten{\longrightarrow}0\quad,
\]
qu'on peut interpréter comme
\diagram{\begin{tikzcd}[ampersand replacement=\&,column sep=small,row sep=small]
(**)\quad 0\arrow[r] \& \underline{t}_A^{\handwritten{\vee}}\arrow[r]\arrow[d,equal] \& \mathbb{D}^{\handwritten{*}}(A)_S\arrow[r] \& \underline{t}_{A^*}\arrow[r] \&0\ \unsure{-|{}}\ .\\
\& \handwritten{\underline{\omega}_A}
\end{tikzcd}}

Le foncteur $(*)$ et la suite exacte $(**)$ sont fonctoriels en $A$ et compatib\deleted{\unsure{l}}les avec image inverse. On peut \deleted{envisager} considérer que le foncteur $(*)$ se factorise en
\[
\mathbb{D}^*:\quad\text{Schémab}(S)\handwritten{\longrightarrow}\mathrm{Crisloclib}_{\mathrm{nilp}}\mathrm{fil}(S)
\]
vers les c\deleted{\unsure{r}}ristaux en modules localement libres ``filtrés'' (en deux crans localement libres). Utilisant cette factorisation, on trouve que pour un objet $S'$ du site cristallin nilpotent de $S$, la\unsure{,|{}} donnée d'un prolongement de $A$ en un schéma \inserted{abélien} $A'$ sur $S'$ revient à la donnée d'un prolongement de la filtration c\deleted{\unsure{e}}anonique $(**\handwritten{)}$ de $D^*(A)_S$ en une filtration (à quotients localement libres\handwritten{)} de $D^*(A)_{S'}$. Plus précisément, on a le

\uline{Théorème de déformation pour les schémas abéliens}. Soient $\deleted{\unsure{E}}S$ un sch\deleted{\unsure{e}}éma, $S'$ un voisinage à puissances div. loc. nilp. de $S$, alors le foncteur naturel
\[
\begin{aligned}
\text{Schémab}(S')\handwritten{\longrightarrow}{}&
\text{catégorie des couples }(A,\deleted{\unsure{D}}\deleted{\unsure{E}}\mathrm{Fil}^1)\text{ d'un schéma abélien}\\
& A\text{ sur }S\text{ et d'un sous-module localement facteur direct}\\
&\mathrm{Fil}^1\ \inserted{\text{de }\mathbb{D}^*(A)_{S'}}\text{ prolongeant }\mathrm{Fil}^1\mathbb{D}^*(A)_S\handwritten{\simeq\underline{\omega}}_A
\end{aligned}
\]
est une équivalence de catégories.

\handwritten{b)} On peut donner une deuxiè\deleted{\unsure{p}}me interprétation du foncteur $\mathbb{D}^*$ sur les

\scanpage{46}
% archive sheet number: 45
schémas abéliens, \inserted{permettant de définir un foncteur qu.inverse du foncteur du théorème.} Pour ceci, on considère pour tout schéma abélien $A$ sur $S$ son e xtension vectorielle universelle, qui est une extension
\[
0\handwritten{\longrightarrow}\underline{t}_{A^*}^{\handwritten{\vee}}\handwritten{\unsure{\mathbin{-}|\longrightarrow}}E(A)\handwritten{\mathbin{-}}A\handwritten{\longrightarrow}0\quad .
\]
Il se trouve qu'on peut trouver un cristal en groupes lisses canonique sur $S$, soit $\underline{E}(\deleted{\unsure{S}}A)$, tel que l'on ait
\[
E(A)\handwritten{\unsure{\simeq|\overset{\sim}{\longrightarrow}}}\underline{E}(A)_S\quad .
\]
La définition est telle que \deleted{pour tout le foncteur}\inserted{$A$} $\underline{E}(A)$ définisse un foncteur
\[
\text{Schémab}(S)\handwritten{\longrightarrow}\text{Cris Groupes lisses}_{\mathrm{nilp}}(S)
\]
compatible aux changements de base, l'isomorphisme précédent étant également fonctoriel et compatible aux changements de base. Il s'ensuit donc que l'on peut définir $\underline{E}(A)_{\deleted{\unsure{S}}}$ comme $S'\handwritten{\longmapsto}\underline{E}(A)_{S'}$, avec
\[
\underline{E}(A)_{S'}=E(A')\quad,
\]
$A'$ désignant un schéma abélien qui prolonge $A$ sur $S'$, \deleted{\unsure{à}} à charge de définir \deleted{un}\inserted{un système transitif d'}isomorphismes canoniques entre les groupes obtenus pour des prolongements différents. Dans le cas où $S$ est d e carac\deleted{\unsure{é}}té\deleted{\unsure{t}}ristique nulle (donc plus question de puissances divisées !) cela est particulièrement simple; on prouve qu'il existe à isomorphisme \uline{unique} près un seul schéma en groupes lisses $E'$ sur l'épaississement $S'$ de $S$ qui prolonge $E$ (NB il suffirait même d'un nilépaississement, comme on voit par passage à la limite à partir du cas noethérien): c'est le résultat déjà signalé dans ma vieille lettre à Tate. Dans le cas général, c'est plus délicat, et peut se faire par ``la méthode de l'exponentielle'', qui sera exposée dans l\deleted{\unsure{t}}e contexte analogue des groupes de BT dans la suite du séminaire. On p\deleted{\unsure{e}}eut aussi définir $\underline{E}(A^*)$ par
\[
\underline{E}(A^*)\handwritten{=}R^1f_{A\mathrm{cris}*}(\underline{G}_{m\ A_{\mathrm{cris}}})\quad ,
\]
où $A_{\mathrm{cris}}$ \deleted{\unsure{c}}désigne le topos cristallin \deleted{\unsure{absi|abso|abs}} nilpotent absolu (relativement à $\underline{Z}$): cette méthode peut également s'adapter au cas des groupes de BT.

Quoi qu'il en soit de la \deleted{\unsure{n}}méthode utilisée pour définir $\underline{E}(A)$, on
